\section{Computational study of accepted adaptivity}\label{sec:study}
We use the same canonical operating primitives throughout the exact comparison: $T=8$, $\beta=0.97$, $q_0=0.5$, $z_0=1$, and $D_z=z+D_0$, where $D_0=0,1,2,3$ has probabilities $0.1,0.3,0.4,0.2$. The regime transition matrix is
\[
 P=\begin{pmatrix}.70&.25&.05\\.15&.70&.15\\.05&.25&.70\end{pmatrix}.
\]
Stock is initially $0,0.5,\ldots,6$ and the historical tier menu is $0,0.1,\ldots,1$. Coefficients are $c=0.3$, $h=0.6$, $p=0.65$, $\kappa=1.2$, $m=1.5$, $\lambda=0.6$, and $r_z=0.6+1.4z$. Customer service valuation is $\nu=2$. A common retainer of 6 can be included on both sides of every provider comparison; the tabulated rewards exclude that constant. The continuous accepted hierarchy permits both stock and tier intervals, using the stock-identification proof rather than an unverified interpolation of the menu optimum.

The initial finite model has 264 nonterminal state nodes and 37,752 stock--tier choices over those nodes. It is small enough for an exact certificate, not presented as a demanding neural benchmark. Its function is to identify economic mechanisms and provide an independently checkable reference. The coupled portfolio experiment has a different function: it tests the accuracy, acceptance, and repeated-solve cost of a learned representation at larger dimension.

\subsection{Accepted information and menu resolution}
Table~\ref{tab:accepted-hierarchy} gives the accepted information decomposition. The time increment is $0.000816876783$, the subsequent regime increment is $0.616558016063$, and the inherited-tier increment is $0.019290378765$. These four optimizations have identical customer, fill, and physical-cost reservations. The continuous full-state optimum has provider reward $-5.233445701435\ldots$ and exactly the static customer's utility and service. Its exact coefficients are available for direct policy deployment without a neural network.

The older unconstrained information comparison is a different experiment: its inherited-tier increment is $0.050593409253$. We retain that result in the companion but do not substitute it for the accepted increment. Likewise, the historical frozen tier of $0.5$ is not used as the principal accepted benchmark. The continuous static optimum improves its reward by $0.076946642417$ before any adaptivity is introduced.

Table~\ref{tab:fine-menu} checks whether coarse tiers and randomization generate the accepted benefit. The randomized optimum on 320 intervals is within approximately $1.10\times10^{-5}$ of the exact continuous optimum. A deterministic incumbent on that menu gains $0.63663320$ over the continuous static contract. The best continuous protocol needs no lottery at all. The extracted deterministic incumbents are not claimed to be globally optimal on their discrete menus; one less favorable extraction at 20 intervals is retained rather than replaced by an invented monotone sequence.
\begin{table}[!ht]
\centering
\caption{Accepted fine-menu optima and deterministic incumbents}\label{tab:fine-menu}
\begin{tabular}{rrrr}\toprule
Tier intervals & Randomized optimum & Deterministic incumbent & Incumbent gap\\
\midrule
10 & -5.244970661 & -5.247006823 & 0.002036162\\
20 & -5.237637537 & -5.284466721 & 0.046829184\\
40 & -5.234637362 & -5.236634206 & 0.001996845\\
80 & -5.233738433 & -5.236860106 & 0.003121673\\
160 & -5.233499396 & -5.233703759 & 0.000204363\\
320 & -5.233456635 & -5.233477772 & 0.000021137\\
\bottomrule\end{tabular}
\par\smallskip\begin{minipage}{.98\textwidth}\small Same customer, fill, and cost reservations as Table~\ref{tab:accepted-hierarchy}. Only the first randomized row also has an independent rational R6 certificate. Other menu bounds are floating-point Lagrangian bounds. Deterministic entries are feasible incumbents, not claimed deterministic optima; their nonmonotone gaps reflect extraction, not nonnested feasible sets. The exact continuous optimum is $-5.233445701435\ldots$.\end{minipage}
\end{table}


\subsection{Robustness under the actual acceptance constraints}
The new sensitivity study solves the accepted full-stock occupation program, not only a provider Bellman relaxation. Each perturbed primitive has a newly optimized continuous static reference. Its fill, customer utility, and physical cost set the reservations unless that row explicitly changes a reservation. When customer utility is raised, we also recompute the accepted static tier before reporting its gain. A comparison with an infeasible old reference would not establish accepted dominance.

The 31-case design varies adjustment cost, maintenance curvature and the maintenance power, the indemnity coefficient, premium intercept and slope, customer service valuation, regime persistence, demand dispersion, demand probabilities, utility and fill reservations, and the stock grid. Tier resolution is varied separately in Table~\ref{tab:fine-menu}. For every feasible LP we record the primal reward, flow residual, constraint slack, reduced-cost residual, and primal--dual gap. These floating-point records are numerical evidence, not rational certificates for every perturbation. Continuous information decompositions are reported only when the stock-support and interior quadratic checks succeed.

Table~\ref{tab:robustness} displays selected cases; EC Table~\ref{tab:robust-all} contains every row. Of 31 accepted programs, 29 are feasible; 28 of those exhibit a positive gain over the applicable continuous accepted static comparator. A zero regime-premium slope removes the accepted gain in the reported instance. Increasing the fill floor by 0.015 is infeasible under the original physical-cost cap, and remains infeasible with the tested cap relaxation of 0.6. Both adverse rows are retained. Thus the conclusion is neither universal feasibility nor positive gain for arbitrary service promises.
\begin{table}[!ht]
\centering
\caption{Accepted-problem sensitivity: selected cases}\label{tab:robustness}
\begin{tabular}{lrrrr}\toprule
Case & Grid gain & Active & Lottery & Memory gain\\
\midrule
canonical & 0.625140 & FUC & 1 & 0.019290\\
lam = 0.1 & 0.728444 & FUC & 1 & 0.001115\\
lam = 4.0 & 0.361126 & FUC & 1 & 0.089863\\
power = 1.5 & 0.983858 & FUC & 1 & ---\\
power = 3.0 & 0.377330 & FUC & 1 & ---\\
persistence = 0.6 & 0.580915 & FUC & 1 & 0.022693\\
dispersion = 1.2 & 0.603774 & FUC & 1 & 0.018567\\
utility delta = 0.2 & 0.585697 & FUC & 1 & 0.018221\\
fill delta = 0.015 & infeasible & --- & --- & ---\\
zero regime premium & 0.000000 & FUC & 0 & ---\\
\bottomrule\end{tabular}
\par\smallskip\begin{minipage}{.98\textwidth}\small Definitions and all 31 cases appear in EC Table~\ref{tab:robust-all}. Reservations and static tiers are recomputed within each perturbed problem. The higher-fill case is infeasible under its stated physical-cost cap. The zero-slope case has no accepted gain.\end{minipage}
\end{table}


Increasing the maintenance power from two to three leaves a positive accepted grid gain of approximately $0.37733$; power $1.5$ yields approximately $0.98386$. Altering dispersion about the same mean also leaves positive gains in both tested directions. These results complement, rather than prove, Theorem~\ref{thm:local-gain}. The theorem identifies a customer-neutral direction under explicit strict inequalities, whereas the numerical study allows stock changes and solves the full constrained problem. Under stock identification, changing $\nu$ alone cancels between the policy and the reservation because filled demand is the same; that invariance is an accounting implication, not new empirical evidence.

The inherited-state increment changes with the commitments and primitives. For example, it is about $0.00111$ when $\lambda=0.1$ and $0.08986$ when $\lambda=4$, even though total accepted grid gain falls in this slice. A cost can reduce the overall benefit of adaptation while making information about the installed tier more valuable. This distinction is obscured by reporting only static-versus-dynamic gains.

\subsection{Commitment and implementation}
The customer-exit program imposes 3,280 continuation constraints on the original eight-review tree. The reported rational protocol satisfies every one, has reward $-5.417032209823\ldots$, and retains a gain of $0.453078763224\ldots$. Its upper-bound gap is below $6.95\times10^{-7}$. This result shows directly what survives when customer commitment is relaxed; it does not assume that an ex ante participation check enforces participation later.

For the coupled portfolio study, all 288 submitted decisions are stored as rational tier vectors, together with their nonnegative payment multipliers, graph matrices, customer slack, and loss bounds. An independent checker uses Python integer sparse multiplication, verifies the graph coercivity condition, and recomputes every fraction without calling an optimizer. At 1,024 services, the maximum learned value--gradient loss bound is $8.0436\times10^{-7}$ per service, compared with a minimum realized gain of approximately $0.016067$ per service over the optimized uniform tier. This is a decision-relevant bound on the accepted update, not the earlier finite-chain global residual budget.

The learned fit has a modest average loss advantage over its equal-oracle value-only counterpart, but the paired seed intervals contain zero. Conjugate gradients have substantially better accuracy in this experiment. The learned critic nevertheless supplies a compact compatible representation, accepted decisions, transfer across graph sizes, and a measured deployment tradeoff relative to refactorization. The result is a reproducible method comparison, not evidence that every numerical control problem requires a neural method.

\section{Operational implications}
An accepted tier revision reallocates payments and maintenance resources; it is not merely a provider's change of performance metric. Three ingredients make its value interpretable. The static alternative must be optimized. Customer service and participation must be enforced in the same comparison. Information classes must be optimized separately under those commitments. With those checks, the canonical gain remains positive, its sources can be separated, and costly inherited-state information has a measurable incremental value.

The exact scalar solution, the interim-participation extension, and the learned portfolio update answer successive implementation questions. They use different computational representations but the same acceptance logic. Complete contractibility remains an assumption, private-information incentive constraints remain distinct, and the portfolio benchmark is a coupled accepted review rather than a solution of an unrestricted high-dimensional multistage problem. The retained continuous-rate and diffusion analyses are labeled as extensions, with their boundary and discretization assumptions intact.

\paragraph{Reproducibility.}
Separate entry-point scripts implement exact information values, robustness, interim participation, portfolio learning, inaction, and independent verification in \path{revisions/or-r7-20260921/}. The scripts record rational certificates separately from floating-point sensitivity results. Exact quantities are compared as fractions; tolerances are used only for the explicitly numerical programs. Historical source directories remain unchanged. The response document maps every referee concern to the current result and records the remaining scope boundaries.
